\begin{table}[!h]
\centering
\caption{\label{tab:sde}Standardized Effect Sizes}
\centering
\begin{threeparttable}
\begin{tabular}[t]{lrrrrrl}
\toprule
Outcome & $\hat{\beta}$ & SE & SD($Y$) & SDE & SE(SDE) & Classification\\
\midrule
Homicide rate (level) & -0.286 & 2.473 & 52.51 & -0.0055 & 0.0471 & Small negative\\
Homicide rate (log) & -0.0454 & 0.1488 & 1.58 & -0.0288 & 0.0943 & Small negative\\
\bottomrule
\end{tabular}
\begin{tablenotes}
\item \textit{Notes:} 
\item textbf{Country:} Brazil. textbf{Research question:} Does an exogenous increase in unconditional fiscal transfers (FPM) to municipal governments affect local homicide rates? textbf{Policy mechanism:} Brazil's Fundo de Participac{c}~{a}o dos Munic'{i}pios distributes federal tax revenue to municipalities through a formula with 17 sharp population thresholds; crossing a threshold increases per-capita transfers by approximately 20%, funding public employment, infrastructure, and services. textbf{Outcome definition:} Homicide rate per 100,000 population (ICD-10 codes X85--Y09 from DATASUS SIM cause-of-death records). textbf{Treatment:} Binary --- municipality annual population above vs. below the nearest FPM threshold (sharp RDD). textbf{Data:} DATASUS SIM mortality data (via IPEADATA) and IBGE population estimates, 2001--2021, municipality-year level, N = 105,767 municipality-years. textbf{Method:} Multi-cutoff sharp RDD (Cattaneo et al. 2016), MSE-optimal bandwidth, triangular kernel, robust standard errors clustered by state. textbf{Sample:} All Brazilian municipalities (5,570) with valid population and mortality data. Classification thresholds: Large ($|$SDE$|$ $>$ 0.15), Moderate (0.05--0.15), Small (0.005--0.05), Null ($<$ 0.005). Classification labels refer to the magnitude of the standardized point estimate, not to statistical significance. ``Null'' denotes a near-zero effect size ($|$SDE$| < 0.005$), not a failure to reject a null hypothesis.
\end{tablenotes}
\end{threeparttable}
\end{table}
